% Draft: §1, Introduction, for the restructured paper (rewrite method agreed
% 17 Oct 2026; see the evidence repo's DECISIONS.md). Drafted 18 Oct from the
% structure review's plan (reviews/structure-muse-20261017-115410/MUSE.md,
% "Section 1" and "The draft opening paragraph"; budget 850 words) and the
% §§2–7 drafts, whose claims and pages it summarizes. The old main file's
% introduction was written on version B (fusion kept for the rich; a
% determinative building an NP in every function) and isn't carried over.
% Pending: Wu (2021), Exocentric Noun Phrases in English (UCL thesis), which
% reportedly argues against fusion as a single solution for NPs without a head
% noun. Not yet read, so not cited; position it in the rivals paragraph once
% read. SECTION-NOTES asks §1 to cover term-of-office once (an ordinary
% modifier on ordinary Head).
%
% Label: sec:intro. Relies on sec:analyses, sec:det-head, sec:mod-head,
% sec:once, sec:costs, sec:fused-relatives.

\section{Introduction}\label{sec:intro}

When one expression seems to have two functions at once, what should persuade us that it does? In \mention{some apples}, \mention{some} determines the noun phrase, and the noun \mention{apples} is its head. In \mention{I'll take \bk{some}}, \mention{some} is the whole noun phrase, and there's no noun to be its head. One influential analysis gives \mention{some} both functions there: it's the determiner of its noun phrase and the phrase's head at once. This paper asks what evidence could justify that doubling.

Two distinctions run through it. A word's category is the set it belongs to, such as noun, determinative, or adjective. Its function is the relation it bears in its phrase, such as head, determiner, or modifier. So \mention{some} is a determinative in both examples, and what's at issue is its functions.

The doubling is called \term{fusion of functions}. Its most explicit statements are in \textit{The Cambridge Grammar of the English Language} (\textit{CGEL}; \citealp{huddleston2002}) and in \textcite{Payne2007}. In \mention{take some}, \textit{CGEL}'s \mention{some} is a fused determiner-head. In \mention{the lucky few} and \mention{the rich}, \mention{few} and \mention{rich} are fused modifier-heads, at once modifier and head (pp.~410--421). Payne et al.\ make the device explicit, extend it to \mention{once}, \mention{twice}, and \mention{thrice}, and say what it costs.

It departs from canonical structure twice. A constituent with two functions \enquote{will in effect have two mothers}, which they call \enquote{the major departure from canonical principles of structure} (p.~570). And the head of a noun phrase can turn out to be a determinative or an adjective, not a noun, \enquote{the second departure} (p.~571). The question is when a grammar should accept the first.

Many syntacticians would sidestep it. The determinative could head the whole noun phrase, as on the DP hypothesis \citep[see][]{millerpullum2022tete,salzmann2020npvsdp}, or a silent noun could follow \mention{some} \citep{lobeck1995ellipsis,kester1996adjectival}. Each route gives up one of three commitments \textit{CGEL} argues for (Section~\ref{sec:analyses}): where a noun phrase has a noun, the noun is its head and the determinative a dependent; \mention{some} keeps its category when it stands alone; and there's no silent head noun. This paper keeps all three, since they're the premises that led to fusion, and asks whether fusion follows from them. It doesn't. The result is conditional on those commitments, but the criterion behind it applies to a second function in any analysis.

The alternative I call \term{ordinary Head}. The phrase that \textit{CGEL} fuses has one function, head of the \term{nominal}, as a common noun does in \mention{take \bk{apples}}. The nominal is the unit between noun and noun phrase, such as \mention{old man} in \mention{the old man}. Since the two analyses describe the same strings, choosing between them takes a criterion, and I put the burden on the second function. A property of form counts as evidence for it if it varies with that function and not with the word alone, and the evidence is weighed against the extra structure (Section~\ref{sec:analyses}).

The test finds no such property where a word that's elsewhere a modifier in a nominal heads the nominal itself, as in \mention{the lucky few}, \mention{the rich}, and \mention{nylon} in \mention{I prefer cotton shirts to \bk{nylon}} (Section~\ref{sec:mod-head}), nor in \mention{once}, \mention{twice}, and \mention{thrice}, on which Payne et al.\ built their case (Section~\ref{sec:once}). Where a determinative stands alone, as in \mention{take some}, it finds one, taken up below, and no other (Section~\ref{sec:det-head}). The genitive, the form \textit{CGEL} ties to a fused determiner-head (p.~420), divides these words one by one, not by function, both standing alone and after a determiner. Payne et al.'s term-of-office \mention{once} and \mention{twice} (\mention{the once mayor}, \mention{the twice president}) come out as ordinary modifiers.

The one property the test finds is the absence of an adjective before a bare determinative (\ungram{\mention{Lucky few escaped}}).\footnote{As in \textit{CGEL}, * marks an unacceptable form, \judgmark{?} a questionable one, \judgmark{!} a non-standard one, and \judgmark{\%} one found only in some varieties.} Fusion derives it, where ordinary Head has to state it as a condition, one of a kind \textit{CGEL} already states for pronouns, so it counts as evidence for the determiner function. It's one fact, though, with an exception in \mention{Poor everybody!}, and Section~\ref{sec:costs} weighs it against the generality ordinary Head gains and the constituents with two parents it avoids. On balance, then, the phrase \textit{CGEL} fuses is an ordinary head, though the evidence doesn't force that conclusion.

Fused relatives are another matter (Section~\ref{sec:fused-relatives}). In \mention{\bk{Whoever told you that} can't have read the report carefully}, \mention{whoever} is head of the subject noun phrase and also the subject of \mention{told you that} \citep[1076]{huddleston2002}. On the alternative, the relative word simply heads the noun phrase and the clause modifies it \citep{bresnan1978free}. But the second function shows in form, twice categorically and once as a preference. Fused relatives relativize their subject, which bare relatives in standard English can't do. They're never infinitival, though relatives without a relative word can be. And the form of \mention{who} and \mention{whoever} tends to answer to the function inside the clause. Fusion gets all three from the rules for relative clauses with a relative phrase, where the head analysis needs licences no other head has. So fusion is favoured there, though not forced, and it's the one construction where this paper keeps it.
